\section{以太网 PHY：从 MAC 帧到介质信号}

网卡驱动建立描述符并启动 DMA 后，MAC 得到的是一串帧字节；线缆或光纤传输的却是经过编码、加扰、纠错和模拟发送的符号。二者之间由协调子层、介质无关接口、PCS、PMA 和 PMD/MDI 等层次连接。不同速率和介质采用的编码与训练细节不同，但从 MAC 到 PHY 的责任边界保持一致：MAC 管理帧和链路层地址，PHY 管理符号、时钟恢复、信道适配与介质电气/光学接口。\cite{ieee8023-baselines}

\topic{发送路径逐层改变表示形式}
发送时，MAC 先形成前导码、帧头、有效载荷和 FCS，并通过 MII 家族接口把数据交给 PHY。PCS 完成编码、加扰、lane 分配与对齐标记；需要前向纠错时，FEC 编码器加入冗余；PMA 负责串并转换、lane 复用和时钟关系；PMD 最终驱动铜缆、背板或光模块。每层都可能扩大数据量或改变时钟域，因此“MAC 每秒交出多少有效载荷”不能直接等同于“介质上需要多少符号带宽”。

接收方向按相反顺序推进。模拟前端先均衡并判决采样，时钟数据恢复电路从信号中恢复采样相位，PMA 形成稳定符号流，FEC 尝试纠错，PCS 解码并恢复 lane 顺序，MAC 最后校验 FCS 并决定接收、统计错误或丢弃。一个 FCS 错误只说明帧到达 MAC 时内容不一致；要继续定位，还需查看纠错计数、符号错误、lane 对齐、信噪比和收发光功率等 PHY 证据。

\begin{pdfdiagram}[x=1cm,y=1cm]
  % 发送路径和接收路径分开占两排，避免返回箭头穿过节点。
  \node[hw label, font=\small\bfseries] at (-0.15,3.85) {发送};
  \node[hw state, minimum width=2.05cm, align=center] (txdma) at (1.05,3.2) {主机内存\\TX 描述符与帧};
  \node[hw control, minimum width=2.05cm, align=center] (txmac) at (3.8,3.2) {MAC\\地址、长度、FCS};
  \node[hw compute, minimum width=2.05cm, align=center] (txpcs) at (6.55,3.2) {PCS / FEC\\编码与加扰};
  \node[hw compute, minimum width=2.05cm, align=center] (txpma) at (9.3,3.2) {PMA / PMD\\串行化与驱动};
  \draw[hw bus] (txdma) -- (txmac);
  \draw[hw bus] (txmac) -- (txpcs);
  \draw[hw bus] (txpcs) -- (txpma);

  \node[hw label, font=\small\bfseries] at (-0.15,0.95) {接收};
  \node[hw state, minimum width=2.05cm, align=center] (rxdma) at (1.05,0.3) {RX completion\\DMA 到主机页};
  \node[hw control, minimum width=2.05cm, align=center] (rxmac) at (3.8,0.3) {MAC\\FCS 与帧过滤};
  \node[hw compute, minimum width=2.05cm, align=center] (rxpcs) at (6.55,0.3) {PCS / FEC\\纠错与解码};
  \node[hw compute, minimum width=2.05cm, align=center] (rxpma) at (9.3,0.3) {PMA / PMD\\采样与时钟恢复};
  \draw[hw return] (rxpma) -- (rxpcs);
  \draw[hw return] (rxpcs) -- (rxmac);
  \draw[hw return] (rxmac) -- (rxdma);

  \node[hw memory, minimum width=2.15cm, minimum height=1.2cm, align=center] (medium) at (11.85,1.75) {介质\\铜缆 / 光纤};
  \draw[hw bus] (txpma.east) -- (10.7,3.2) -- (10.7,2.25) -- (medium.north west);
  \draw[hw return] (medium.south west) -- (10.7,1.25) -- (10.7,0.3) -- (rxpma.east);
\end{pdfdiagram}
\diagramnote{一帧以不同表示穿过以太网接口。主机内存保存字节和描述符，MAC 处理帧语义，PCS/FEC 处理码块与纠错，PMA/PMD 处理 lane、采样和介质信号；接收完成只有在 PHY 恢复数据、MAC 接受帧并把内容 DMA 到主机页后才成立。每层都留下不同的错误证据。}

\topic{1 Gb/s 数值案例：64 B 最小帧并不只占 512 ns}

以不带 VLAN 标签的经典以太网帧为例，从目的 MAC 地址到 FCS 的最小帧长为 64 B，其中有效载荷不足时要填充到 46 B。介质上还要发送 7 B 前导码和 1 B SFD，并在下一帧前留出相当于 12 B 的帧间隔。因此，一次最小帧占用的线速时间为
\[
(8+64+12)\times8\ \mathrm{bit}\div10^9\ \mathrm{bit/s}
=672\ \mathrm{ns}.
\]
一个 1518 B 普通最大帧对应
\[
(8+1518+12)\times8\ \mathrm{bit}\div10^9\ \mathrm{bit/s}
=12.304\ \mu\mathrm{s}.
\]
这两个数字解释了为什么“每秒帧数”和“每秒有效载荷字节数”不是同一个指标：短帧中前导码、帧头、FCS 与帧间隔所占比例更大；驱动、DMA 和中断还会增加不在线路字节数中的主机开销。\cite{ieee8023-baselines}

\begin{longtable}{L{0.18\linewidth}L{0.23\linewidth}L{0.29\linewidth}L{0.20\linewidth}}
\toprule
边界 & 表示形式 & 本层确认什么 & 本层不能确认什么 \\
\midrule
TX 描述符 & 地址、长度、offload 与所有权位 & 驱动已经发布一块可 DMA 缓冲区 & 帧已经离开 PHY \\\tablerowrule
MAC & 前导、地址、类型/长度、payload、FCS & 帧格式成立并完成 FCS 生成或检查 & 模拟眼图和 FEC 裕量 \\\tablerowrule
PCS / FEC & 码块、加扰数据、对齐标记与冗余 & 码块同步，错误在能力范围内已纠正 & 远端应用已经消费数据 \\\tablerowrule
PMA / PMD & 串行符号、采样相位与电气/光学幅度 & lane 锁定并满足当前接收判决 & MAC 帧地址与上层校验正确 \\\tablerowrule
RX completion & 缓冲区、长度、状态与校验结果 & NIC 已把被接受的帧写入主机内存 & 应用线程已经运行并读取它 \\
\bottomrule
\end{longtable}

\topic{自动协商先确定共同能力}
链路两端上电后先交换能力，选择双方共同支持的速率、双工方式以及适用的流控或 FEC 参数。某些高速铜缆和背板链路还要执行链路训练：发送端提供规定序列，接收端测量信道响应并请求调整发射均衡，直到误码和裕量满足进入数据态的条件。协商完成只表示参数达成一致；训练、块锁定、lane 对齐和错误率稳定后，链路才真正可承载 MAC 流量。

排障时应把状态分成三层：没有协商结果，优先检查介质、对端能力和管理配置；已经协商但训练失败，检查损耗、端接、均衡和 lane；链路已进入数据态但错误持续增长，再检查 FEC 纠错余量、温度、连接器与串扰。直接重启驱动可能暂时重新训练，却不能证明物理问题消失。

\topic{FEC 用确定延迟换取误码容限}
假设接收端在一个 FEC 码块内检测到少量错误，译码器可利用冗余恢复原始数据并增加“已纠正码字”计数；错误超过纠错能力时，只能报告不可纠正码字，上层帧校验通常也会失败。FEC 改善的是给定信道条件下的有效误码率，不会增加模拟眼图本身的裕量，并且会引入编码冗余和固定流水延迟。

因此健康链路也可能持续出现少量已纠正错误。判断是否退化要观察单位时间纠错率、不可纠正错误、链路重训练次数和环境变化，而不是看到非零计数就立即更换硬件。若纠错率随温度或线缆弯折显著上升，它就是物理裕量收窄的早期证据。

\topic{从收帧异常回溯到物理层}
当应用报告吞吐下降时，先确认 MAC 是否出现丢帧、暂停帧或环耗尽；若 MAC 错误不多，再检查驱动队列和 CPU 调度。若 FCS、符号或 FEC 计数同步上升，则沿 PCS、PMA、PMD 回溯。这个顺序能避免把物理链路误判成软件队列问题，也能避免在 RX 环已经耗尽时反复调整模拟均衡。完整证据链应同时保存协商结果、训练状态、PHY 计数器、MAC 计数器、驱动队列和应用时间线。

\topic{案例：吞吐下降但 link 仍为 up}

假设应用吞吐从 940 Mb/s 降到 300 Mb/s，而管理界面仍显示 1 Gb/s、全双工。第一种可能是主机侧 RX 环耗尽：PHY 与 MAC 计数正常，但 \code{rx\_no\_buffer}、丢包和 CPU 软中断积压上升，此时应增加队列处理能力或修正中断亲和性。第二种可能是介质裕量下降：FCS、符号错误或已纠正 FEC 码字随时间增加，并伴随重训练，此时应检查线缆、连接器、光功率、温度和均衡。第三种可能是流控：收到大量 pause 帧时，发送端会按协议停顿，链路不掉线却无法持续发满。

定位的关键不是先重启驱动，而是把同一时间窗内的应用速率、RX/TX 环水位、MAC FCS 计数、PHY 符号/FEC 计数、pause 帧和重训练次数对齐。队列水位先恶化而 PHY 计数稳定，根因更接近主机；PHY 错误先增加而队列随后空闲，根因更接近介质。只有最早变化的证据才能指向最早失效边界。
